\chapter{Research Plan (Kế hoạch nghiên cứu)}

\section{Kế hoạch 3 phase}

Project chia thành 3 phase.
Mỗi phase tạo ra đầu vào cho phase sau: phase 1 chốt hướng đi, phase 2 tạo bằng chứng, phase 3 biến bằng chứng thành sản phẩm.

\begin{table}[H]
  \centering
  \caption{Tổng quan 3 phase}
  \begin{tabularx}{\linewidth}{>{\raggedright\arraybackslash}p{3.3cm} X X}
    \toprule
    \textbf{Phase} & \textbf{Mục tiêu} & \textbf{Kết quả chính} \\
    \midrule
    1. Hiểu bài toán & Cả nhóm có chung bức tranh tổng thể & Dataset và nguồn pose đã chốt, danh sách approach, kế hoạch phân công \\
    \addlinespace
    2. Code pipeline của mỗi approach, kèm chạy thử & Có số liệu thật để so sánh & Demo và bảng benchmark của 4 approach \\
    \addlinespace
    3. Xây tool và bàn giao & Gom kết quả thành sản phẩm dùng được & Tool, technical report, bàn giao cho mentor \\
    \bottomrule
  \end{tabularx}
\end{table}

\subsection{Phase 1: Literature review và thống nhất hướng đi}

Mục tiêu của phase này là để cả nhóm hiểu cùng một bài toán trước khi chia việc.

Công việc:
\begin{itemize}
  \item Đọc tài liệu và tổng hợp các phương pháp, approach và pipeline hiện có.
  \item Thảo luận nhóm để thống nhất bức tranh tổng thể: bài toán, dataset, phương pháp, các thành phần cần có và các công việc cần làm.
  \item Khảo sát dataset cho hai track pinhole/fisheye + LiDAR, kiểm tra dữ liệu có đủ calibration, pose và timestamp không.
  \item Chốt nguồn pose: dùng pose có sẵn của dataset, hay phải chạy ước lượng pose.
  \item Phác thảo high-level architecture và các module contract ở mức đủ để chia việc.
  \item Thống nhất phần dùng chung: format dữ liệu, hệ tọa độ, cách đo và cách ghi kết quả.
\end{itemize}

Kết quả cần có:
\begin{itemize}
  \item Dataset và nguồn pose đã chốt.
  \item Danh sách 4 approach được chọn và lý do chọn.
  \item Bản nháp architecture và evaluation protocol.
  \item Phân công người phụ trách từng approach.
\end{itemize}

\subsection{Phase 2: Chạy thử 4 approach và benchmark}

Nhóm chia nhau tìm hiểu 4 approach/pipeline/method tốt nhất.
Mỗi approach phải được chạy thật, có demo và số liệu, không dừng ở việc đọc paper.

Mỗi người phụ trách một approach và làm theo các bước:
\begin{enumerate}
  \item Đọc paper và code, ghi lại input, output, các module trong pipeline và yêu cầu phần cứng.
  \item Dựng môi trường và chạy thử trên dữ liệu mẫu của repository để xác nhận code hoạt động.
  \item Chuyển dataset đã chốt sang định dạng mà approach yêu cầu.
  \item Chạy thử một đoạn ngắn từ đầu đến cuối: dữ liệu $\to$ calibration/pose $\to$ reconstruction $\to$ render ảnh giữ lại $\to$ tính số đo.
  \item Chạy trên toàn bộ đoạn dữ liệu đã thống nhất và tạo demo.
  \item Ghi lại kết quả theo mẫu chung, kể cả các trường hợp thất bại.
  \item Hoàn thành bảng ánh xạ module và adapter tối thiểu, để approach sớm chạy được trên tool.
\end{enumerate}

Mọi approach được đánh giá trên cùng dữ liệu và cùng cách đo:
\begin{table}[H]
  \centering
  \caption{Nội dung evaluation/benchmark cho mỗi approach}
  \begin{tabularx}{\linewidth}{>{\raggedright\arraybackslash}p{3.5cm} X}
    \toprule
    \textbf{Nhóm} & \textbf{Cần ghi lại} \\
    \midrule
    Chất lượng hình ảnh & Render lại các ảnh không dùng để train (held-out views), so với ảnh thật, ví dụ bằng PSNR \\
    \addlinespace
    Chi phí & Thời gian chạy, VRAM, phần cứng cần dùng, lượng thao tác thủ công \\
    \addlinespace
    Khả năng hỗ trợ & Pinhole, fisheye, LiDAR và những phần approach không làm được \\
    \addlinespace
    Khả năng tích hợp & Module thay thế được, format output, khả năng đưa vào CARLA \\
    \addlinespace
    Failure cases & Trường hợp lỗi, nguyên nhân dự đoán, cách khắc phục đã thử \\
    \bottomrule
  \end{tabularx}
\end{table}

Song song, có thể chạy CARLA với bản đồ có sẵn để kiểm tra phần mô phỏng độc lập với reconstruction.
Nếu kết quả thực tế không khớp với thiết kế ban đầu, cập nhật lại architecture và contract.

Kết quả cần có:
\begin{itemize}
  \item Demo và log chạy của từng approach.
  \item Bảng benchmark so sánh 4 approach.
  \item Nhận xét approach nào phù hợp với track nào, kèm bằng chứng.
  \item Danh sách module và contract đã chỉnh theo kết quả thực tế.
\end{itemize}

\subsection{Phase 3: Xây dựng tool, hoàn thiện và bàn giao}

Phase này tổng hợp kết quả của phase 2 thành một tool hoàn chỉnh.

Công việc:
\begin{itemize}
  \item Chuẩn hóa các module đã chạy được của 4 approach theo contract chung.
  \item Xây dựng tool để chọn method trong từng module và đổi thứ tự module theo từng approach.
  \item Chạy thử các pipeline trên tool, so kết quả với số liệu ở phase 2 để phát hiện sai lệch.
  \item Sửa lỗi và cải thiện tool.
  \item Kiểm tra CARLA integration feasibility và chạy lại để kiểm tra tính nhất quán.
  \item Hoàn thành technical report.
  \item Bàn giao code, hướng dẫn sử dụng, demo và báo cáo cho mentor đánh giá.
\end{itemize}

Kết quả cần có:
\begin{itemize}
  \item Tool chạy được các pipeline đã chọn.
  \item Scene demo và bảng đánh giá cuối cùng.
  \item Technical report và bộ tài liệu bàn giao.
\end{itemize}

\begin{warning}
Các mục sau chưa được xác nhận là deliverable bắt buộc: chạy toàn bộ Autoware và hoàn thành neural sensor closed-loop.
Cần hỏi mentor trước khi đưa vào phạm vi.
\end{warning}

\section{Liên hệ giữa phase và tài liệu}

\begin{table}[H]
  \centering
  \caption{Phase và phần tài liệu tương ứng}
  \begin{tabularx}{\linewidth}{l l X}
    \toprule
    \textbf{Phase} & \textbf{Thời gian đề xuất} & \textbf{Tài liệu} \\
    \midrule
    Phase 1 & Tuần 1 & Part I: chốt dataset, protocol, 4 approach; các chương khác ở mức đủ dùng và cập nhật dần \\
    Phase 2 & Tuần 2--4 & Part II: mỗi người hoàn thành chương của mình \\
    Phase 3 & Tuần 5--6 & Part III: tổng hợp, tích hợp tool, đánh giá CARLA feasibility \\
    \bottomrule
  \end{tabularx}
\end{table}

\section{Responsibilities}

\begin{table}[H]
  \centering
  \caption{Phân công}
  \begin{tabularx}{\linewidth}{X l}
    \toprule
    \textbf{Công việc} & \textbf{Người phụ trách} \\
    \midrule
    Approach A & TBD \\
    Approach B & TBD \\
    Approach C & TBD \\
    Approach D & TBD \\
    Tool và script đánh giá dùng chung & TBD \\
    Road logic, traffic và CARLA & TBD \\
    \bottomrule
  \end{tabularx}
\end{table}

\section{Dependencies}

\begin{itemize}
  \item Chương 4 phải chốt trước khi benchmark, vì mọi số liệu đều phụ thuộc vào protocol.
  \item Dataset và nguồn pose phải chốt trước khi chạy approach.
  \item Script đánh giá dùng chung phải có trước khi so sánh kết quả.
  \item Chạy approach và chạy baseline CARLA không phụ thuộc nhau, nên có thể làm song song.
  \item Tích hợp tool phụ thuộc vào bảng ánh xạ module và adapter của từng approach.
\end{itemize}

\section{Milestones}

Các mốc dưới đây là đề xuất, dựa trên thời gian phase ở mục trên.

\begin{table}[H]
  \centering
  \caption{Mốc kiểm tra}
  \begin{tabularx}{\linewidth}{l X l}
    \toprule
    \textbf{Tuần} & \textbf{Mốc} & \textbf{Mức} \\
    \midrule
    1 & Part I được thống nhất, dataset, protocol và 4 approach đã chốt & M0 \\
    2 & Mỗi approach chạy được trên dữ liệu mẫu, script đánh giá dùng được & M1 \\
    3 & Vertical slice chạy được trên dataset đã chốt & M1 \\
    4 & Có demo và benchmark của 4 approach, Part II hoàn thành & M2 \\
    5 & Các approach chạy được trên tool, bảng so sánh hoàn thành & M2 \\
    6 & Tool ổn định, CARLA feasibility và report hoàn thành, bàn giao mentor & M2, một phần M3 \\
    \bottomrule
  \end{tabularx}
\end{table}

\section{Rủi ro và phương án dự phòng}

\begin{table}[H]
  \centering
  \small
  \caption{Rủi ro chính}
  \begin{tabularx}{\linewidth}{>{\raggedright\arraybackslash}X >{\raggedright\arraybackslash}X >{\raggedright\arraybackslash}X}
    \toprule
    \textbf{Rủi ro} & \textbf{Dấu hiệu sớm} & \textbf{Phương án dự phòng} \\
    \midrule
    Approach không chạy được hoặc thiếu code & Chưa chạy được dữ liệu mẫu sau tuần 2 & Thay bằng approach dự phòng trong shortlist \\
    \addlinespace
    Dataset thiếu calibration hoặc pose & LiDAR chiếu lên ảnh bị lệch & Đổi sequence, đổi dataset hoặc chạy ước lượng pose \\
    \addlinespace
    Thiếu GPU hoặc VRAM & Train quá lâu hoặc hết bộ nhớ & Dùng đoạn dữ liệu ngắn hơn hoặc thuê GPU \\
    \addlinespace
    Khó nối renderer với CARLA & Chưa có phần nối sẵn cho method & Dừng ở baseline CARLA, ghi lại trong feasibility study \\
    \addlinespace
    Trích xuất road logic không kịp & Bản đồ tự động không hợp lệ & Vẽ tay cho một khu vực nhỏ \\
    \addlinespace
    Tích hợp tool bị trễ & Chưa có bảng ánh xạ module cuối tuần 4 & Ưu tiên tool chạy được 2 approach tốt nhất \\
    \bottomrule
  \end{tabularx}
\end{table}

\section{Nhật ký quyết định}

Mỗi quyết định được ghi lại cùng bằng chứng, để biết khi nào cần xem xét lại.

\begin{table}[H]
  \centering
  \caption{Mẫu nhật ký quyết định}
  \begin{tabularx}{\linewidth}{l X X l}
    \toprule
    \textbf{Ngày} & \textbf{Quyết định} & \textbf{Bằng chứng} & \textbf{Trạng thái} \\
    \midrule
     &  &  &  \\
    \bottomrule
  \end{tabularx}
\end{table}
